\begin{abstract}
The activated-range Gram from the preceding fixed-shift analysis is exact but
abstract: its smallest eigenvalue depends on the actual row columns after
residual grouping. We identify their literal support geometry. On one declared
source system, a fixed physical row contributes at most one terminal atom to a
given output fiber. Hence every active normalized row column has diagonal
energy exactly one; all loss is cross-row collision. Summing the factor-exchange
and CRT geometry of an earlier local analysis yields a new uniform row-pair
codegree law. If the root cofactor is in \([R,2R]\), two distinct rows share at
most \(X^{o(1)}(1+J/R)\) reduced output cells, and therefore only \(X^{o(1)}\)
cells on the critical block \(R\asymp J\), even though one output fiber can
still have polynomial row degree. We then connect this sparse incidence to
the lower spectrum. Private-output energy gives a sharp lower frame; an
optimal phase-blind comparison matrix gives weighted diagonal dominance; and
forest Gram graphs admit an exact formula because all edge phases can be
gauged away. Unique-matching minors give phase-uniform injectivity, whereas
ordinary Hall matching does not. A cyclic model with identical support,
degrees, codegrees, and magnitudes is singular or nonsingular according only
to one holonomy phase. Thus low pair codegree is a genuine structural advance,
but it reaches a polynomial physical lower bound only after row-energy
diffuseness or a global expansion/holonomy estimate is proved. No such
endpoint estimate, parity improvement, or prime-pair conclusion is claimed.
\end{abstract}
